\subsection{Calorie Estimation Module}

Pakistan currently lacks an up-to-date national food composition table for
cooked composite dishes~\cite{soomro2024matiari}. Rather than training a black-box
regressor to guess calories from pixels, Nutrify maps each recognised dish
through \texttt{src/calories.py} and \texttt{src/calories.json} to a
deterministic per-serving calorie range (low, typical, and high), which is
scaled linearly by the user's selected portion multiplier.

\paragraph{How each range was obtained.}
Because standard food composition tables list raw ingredients rather than
prepared Pakistani dishes, we constructed recipe-based estimates for one
standard serving of each of the 14 dishes. Primary ingredient energies are
taken from Table~1 of the Indian Food Composition Tables 2017 (IFCT
2017)~\cite{longvah2017ifct}, converted from kilojoules to kilocalories using
$1\text{ kJ} = 0.239\text{ kcal}$ (IFCT computes metabolisable energy using
Atwater factors, General Information, pp.~xxi--xxii). Cooking fats, dairy, and
other staple items not listed in the same form in IFCT (sunflower oil T508 at
$899\text{ kcal/100 g}$, butter ghee T506 at $878\text{ kcal/100 g}$, whole-milk
plain yogurt L520 at $79\text{ kcal/100 g}$, white sugar I502 at
$394\text{ kcal/100 g}$, gram flour A504 at $353\text{ kcal/100 g}$, and fresh
single cream L519 at $193\text{ kcal/100 g}$) are taken from McCance and
Widdowson's Composition of Foods Integrated Dataset (CoFID)~\cite{cofid2021},
as compiled in the open-access Indian Nutrient Databank
(INDB)~\cite{vijayakumar2024indb}.

For each dish, the \emph{typical} value sums the edible ingredient energies for
a home-style or standard restaurant serving, rounded to the nearest
$5\text{ kcal}$. Spices, salt, water, and light garnishes contribute negligible
energy and are omitted; cooking changes moisture content rather than total
energy, and frying oil is counted as the estimated amount absorbed or eaten
rather than the full pan volume. The \emph{low} and \emph{high} values
represent a leaner recipe version (less oil or ghee and leaner meat cuts, such
as skinless chicken breast N003 at $168\text{ kcal/100 g}$ or goat legs O003 at
$160\text{ kcal/100 g}$) and a richer recipe version (more oil or ghee, fattier
cuts such as chicken thigh N002 at $200\text{ kcal/100 g}$ or beef shoulder
O025 at $212\text{ kcal/100 g}$, and a larger main-ingredient portion).
Table~\ref{tab:calories} lists the 14 dishes, their standard serving
descriptions, and their low, typical, and high kilocalorie values from
\texttt{src/calories.json}.

\begin{table}[H]
  \centering
  \small
  \begin{tabular}{|p{4.3cm}|p{4.8cm}|r|r|r|}
    \hline
    \textbf{Dish} & \textbf{Standard serving} & \textbf{Low} & \textbf{Typical} & \textbf{High} \\
    \hline
    Aloo Gosht (\texttt{aloo\_gosht}) & 1 bowl (about 300 g) & 355 & 540 & 700 \\
    Chicken Biryani (\texttt{chicken\_biryani}) & 1 plate (about 350 g) & 610 & 845 & 1080 \\
    Chicken Karahi (\texttt{chicken\_karahi}) & 1 bowl (about 250 g) & 385 & 580 & 770 \\
    Chicken Tikka (\texttt{chicken\_tikka}) & 4 pieces (about 150 g cooked) & 320 & 400 & 515 \\
    Daal (\texttt{daal}) & 1 bowl (about 250 g) & 195 & 295 & 405 \\
    Haleem (\texttt{haleem}) & 1 bowl (about 300 g) & 365 & 600 & 835 \\
    Halwa Puri (\texttt{halwa\_puri}) & 2 puri, 1 small bowl halwa, a little chana (about 250 g) & 620 & 810 & 1085 \\
    Nihari (\texttt{nihari}) & 1 bowl (about 300 g) & 340 & 585 & 800 \\
    Pakora (\texttt{pakora}) & 6 pieces (about 100 g) & 210 & 295 & 405 \\
    Palak Paneer (\texttt{palak\_paneer}) & 1 bowl (about 250 g) & 235 & 355 & 495 \\
    Paratha Roll (\texttt{paratha\_roll}) & 1 roll (about 200 g) & 410 & 545 & 710 \\
    Pulao (\texttt{pulao}) & 1 plate (about 300 g) & 420 & 560 & 715 \\
    Samosa (\texttt{samosa}) & 2 samosas (about 130 g) & 305 & 405 & 510 \\
    Seekh Kebab (\texttt{seekh\_kebab}) & 2 kebabs (about 120 g cooked) & 195 & 315 & 400 \\
    \hline
  \end{tabular}
  \caption{Per-serving calorie estimates in kcal for the 14 dishes, from \texttt{src/calories.json}.}
  \label{tab:calories}
\end{table}

\subsection{Confidence Threshold and Non-Dish Check}

Because a 14-way softmax layer always assigns probabilities that sum to $1$,
the classifier will output a dish name even when given a blurry photo, an
ambiguous dish, or an image that does not contain food at all. To prevent the
application from presenting low-confidence guesses as definite answers,
\texttt{src/confidence\_analysis.py} selects a rejection threshold $\tau$ on
the 640-image validation set. When the top softmax probability falls below
$\tau$, the application withholds the dish decision and the calorie estimate. (In the
final application such a photo is reported as out of scope, and a second
feature-similarity check was added; see the application section above.)

\paragraph{Validation sweep and test evaluation.}
We swept $\tau$ from $0.30$ to $0.90$ in steps of $0.05$ on the validation set
(\texttt{results/confidence\_val.csv}) and applied the selection rule: choose
the smallest threshold achieving at least $0.90$ accuracy on accepted images
while keeping coverage at or above $0.50$. Table~\ref{tab:confidence-val}
shows the full validation sweep. The smallest qualifying threshold is
$\tau = 0.45$ (\texttt{results/confidence\_threshold.json}), which achieves
$0.9173$ validation accuracy at $0.8313$ coverage. Applied once to the 641
held-out test predictions (\texttt{results/confidence\_test.csv}),
$\tau = 0.45$ yields a test coverage of \textbf{0.8456} and raises test
accuracy on accepted images from $0.7988$ (unthresholded) to \textbf{0.8727}.

\begin{table}[H]
  \centering
  \small
  \begin{tabular}{|r|r|r|}
    \hline
    \textbf{Threshold ($\tau$)} & \textbf{Val.\ coverage} & \textbf{Val.\ accuracy (accepted)} \\
    \hline
    0.30 & 0.9375 & 0.8717 \\
    0.35 & 0.9078 & 0.8847 \\
    0.40 & 0.8750 & 0.8964 \\
    \textbf{0.45} & \textbf{0.8313} & \textbf{0.9173} \\
    0.50 & 0.7828 & 0.9301 \\
    0.55 & 0.7469 & 0.9393 \\
    0.60 & 0.7109 & 0.9560 \\
    0.65 & 0.6781 & 0.9608 \\
    0.70 & 0.6312 & 0.9653 \\
    0.75 & 0.5828 & 0.9705 \\
    0.80 & 0.5078 & 0.9815 \\
    0.85 & 0.4250 & 0.9890 \\
    0.90 & 0.3172 & 0.9951 \\
    \hline
  \end{tabular}
  \caption{Validation coverage and accuracy on accepted images across candidate
  thresholds (\texttt{results/confidence\_val.csv}). Bold marks the selected
  threshold $\tau = 0.45$, which achieves $0.8456$ coverage and $0.8727$
  accuracy on the test set (\texttt{results/confidence\_test.csv}).}
  \label{tab:confidence-val}
\end{table}

\paragraph{Informal non-dish check.}
To see how the threshold behaves on out-of-distribution inputs, we ran an
informal check on 6 non-dish photos (\texttt{results/non\_dish\_check.csv}),
shown in Table~\ref{tab:non-dish}. Because there are only 6 images, this is a
qualitative sanity check rather than a statistical benchmark. Four of the 6
images had top softmax confidence between $0.25$ and $0.35$ and were correctly
rejected as \texttt{Not sure}. The remaining two images (a selfie in a grey
T-shirt with a phone at $0.45$ and a selfie in a grey T-shirt at $0.58$) met or
exceeded the $0.45$ threshold and were wrongly accepted as \texttt{Samosa}.

\begin{table}[H]
  \centering
  \small
  \begin{tabular}{|l|r|l|}
    \hline
    \textbf{Image description} & \textbf{Top confidence} & \textbf{App answer ($\tau = 0.45$)} \\
    \hline
    face photo blue shirt & 0.30 & Not sure \\
    Nike logo & 0.35 & Not sure \\
    screen with code & 0.25 & Not sure \\
    selfie grey T-shirt with phone & 0.45 & Samosa \\
    selfie grey T-shirt & 0.58 & Samosa \\
    selfie flexing & 0.33 & Not sure \\
    \hline
  \end{tabular}
  \caption{Informal check on 6 non-dish images (\texttt{results/non\_dish\_check.csv}).}
  \label{tab:non-dish}
\end{table}

\subsection{Calorie and Threshold Limitations}

Several methodological limitations apply to the calorie estimates and the
confidence filter:
\begin{itemize}
  \item \textbf{Recipe estimates, not measurements:} All values in
        \texttt{src/calories.json} are recipe-based calculations from food
        composition tables, not bomb-calorimetry measurements of the photographed
        meal, and they are for general informational use rather than medical or
        clinical dietary advice.
  \item \textbf{Assumed serving sizes and unmeasured portions:} The standard
        serving weight and every ingredient weight in a recipe are fixed
        assumptions. The system does not detect food volume, depth, or plate
        scale from the photo; portion adjustments rely entirely on the user's
        chosen multiplier.
  \item \textbf{Lean and rich recipes, not statistical intervals:} The
        \texttt{kcal\_low} and \texttt{kcal\_high} bounds represent a lean recipe
        variant (less oil or ghee, leaner meat cut) and a rich recipe variant
        (more oil or ghee, fattier cut, larger main-ingredient weight), not
        statistical confidence intervals or standard deviations across kitchens.
  \item \textbf{Cooked poultry weights in IFCT 2017:} IFCT 2017 reports nutrient
        values for raw edible foods \emph{except} for poultry and egg (General
        Information, p.~xx). Accordingly, the chicken weights in our recipes
        (\texttt{chicken\_biryani}, \texttt{chicken\_karahi},
        \texttt{chicken\_tikka}, \texttt{paratha\_roll}, and \texttt{pulao})
        refer to cooked boneless skinless chicken breast (N003,
        $168\text{ kcal/100 g}$) or thigh (N002, $200\text{ kcal/100 g}$),
        whereas red meat, grains, pulses, and vegetables are entered as raw
        edible weights.
  \item \textbf{Omission of IFCT ``chicken, leg, skinless'':} The IFCT 2017
        entry for skinless chicken leg reports an energy of about
        $384\text{ kcal/100 g}$, which is inconsistent with its own reported
        protein and fat content under Atwater conversion; we therefore excluded
        that entry and used only breast (N003) and thigh (N002).
  \item \textbf{Why INDB cooked-dish totals were used only as a cross-check:}
        The Indian Nutrient Databank~\cite{vijayakumar2024indb} includes
        calculated recipes for several Indian dishes (for example Mutton biryani
        ASC122 at $396\text{ kcal}$ per plate, Dal makhani OSR139 at
        $262\text{ kcal}$ per bowl, Spinach paneer ASC215 at $230\text{ kcal}$
        per bowl, and Plain pulao ASC114 at $432\text{ kcal}$ per plate). We did
        not adopt INDB's cooked-dish totals directly because INDB counts all
        frying oil in a recipe as consumed and does not apply cooking yield
        factors; we used INDB only to source UK CoFID~\cite{cofid2021}
        ingredient values and as a sanity cross-check.
  \item \textbf{Closed-set classifier and imperfect rejection:} Because the
        network is trained strictly on 14 classes, it always picks one of those
        14 dishes. The $0.45$ confidence threshold catches most uncertain or
        non-dish photos (4 of 6 in our informal check), but out-of-distribution
        images can still produce a softmax peak at or above $0.45$ and be
        accepted as one of the 14 dishes.
\end{itemize}

